\begin{abstract}
Cross-domain class-incremental learning (CIL) for multimodal remote sensing remains unexplored: a model trained on one geographic domain must continuously learn new land-cover classes from different domains without storing past data. We address this challenge with two contributions. First, we establish a \emph{Marathon CIL} benchmark spanning three HSI+LiDAR datasets (Trento, Houston, MUUFL) across 9 incremental tasks and 32 classes, providing the first standardized cross-domain exemplar-free CIL protocol for multimodal remote sensing. Second, we propose \methodname{}, a modality-asymmetric parameter-efficient fine-tuning strategy grounded in an empirical observation: spectral features exhibit significantly lower cross-domain drift than spatial features (drift ratio $\approx$1:2.3:2.5 for spectral, HSI-spatial, and LiDAR-spatial branches). \methodname{} freezes the stable spectral branch as a semantic anchor after a brief warm-up phase, and adapts the spatial branches with per-task low-rank adapters whose capacity is proportional to each modality's measured drift. On the \backbonename{} backbone, \methodname{} achieves 75.7\% average task-agnostic accuracy, outperforming all exemplar-free baselines including EWC (65.9\%), LwF (68.5\%), and DCPRN (74.6\%), while using only 36K additional parameters. Systematic ablation across 3 domain orderings, multiple seeds, and 12 architectural variants confirms the method's robustness and identifies that asymmetric rank allocation and selective freezing are the key contributors to performance.
\end{abstract}
